\documentclass[12pt]{article}

\usepackage{sbc-template}
\usepackage{graphicx,url}
\usepackage[utf8]{inputenc}
\usepackage[brazil]{babel}
%\usepackage[latin1]{inputenc}  

\begin{document} 

\begin{titlepage}
\begin{center}
    \textbf{\MakeUppercase{Educação Financeira Assistida por IA: Plataforma de Apoio Personalizado com Inteligência Artificial Generativa}}
\end{center}
\vspace{1.5cm}

\noindent\textbf{Identificação dos autores}

\vspace{0.5cm}

\noindent\textbf{1º Autor}

\noindent\textbf{Nome:} João Gabriel Barreto de Araújo Falcão

\noindent\textbf{E-mail:} joao.falcao@cs.unipe.edu.br

\noindent\textbf{Telefone:} (83) 98888-0507

\vspace{0.5cm}

\noindent\textbf{2º Autor}

\noindent\textbf{Nome:} Júlio César Carvalho Santos

\noindent\textbf{E-mail:} julio.santos05@cs.unipe.edu.br

\noindent\textbf{Telefone:} (83) 98722-0156

\vspace{0.5cm}

\noindent\textbf{3º Autor}

\noindent\textbf{Nome:} Perilo de Oliveira Viana Sobrinho

\noindent\textbf{E-mail:} perilo.sobrinho@cs.unipe.edu.br

\noindent\textbf{Telefone:} (83) 98647-3303

\vspace{0.5cm}

\noindent\textbf{4º Autor (orientador)}

\noindent\textbf{Nome:} Ricardo Dantas Silva

\noindent\textbf{E-mail:} ricardods@unipe.edu.br

\noindent\textbf{Telefone:} 

\end{titlepage}
     
\sloppy

\title{EDUCAÇÃO FINANCEIRA ASSISTIDA POR IA: PLATAFORMA DE APOIO PERSONALIZADO COM INTELIGÊNCIA ARTIFICIAL GENERATIVA}

\author{João G. B. A. Falcão\inst{1}, Júlio C. C. Santos\inst{1}, Perilo O. V. Sobrinho \inst{1}, Ricardo Dantas\inst{1} }

\address{Ciência da Computação – Centro Universitário de João Pessoa (UNIPÊ) \\
João Pessoa – PB – Brasil
  \email{\{joao.falcao, julioccsantos, perilo.sobrinho, ricardo.dantas\}@cs.unipe.edu.br}
}

% \begin{document} 

\maketitle

\begin{resumo} 
  No que diz respeito às finanças, os brasileiros enfrentam um desafio significativo, refletido na carência de educação financeira. Pesquisas realizadas pelo Instituto de Pesquisas Sociais, Políticas e Econômicas (IPESPE) sugerem que grande parte dos brasileiros possui pouco ou nenhum 
  entendimento sobre o tema. Embora sistemas tradicionais de controle financeiro auxiliem no acompanhamento de receitas e despesas, a ausência de conhecimento adequado limita a tomada de decisões eficazes. Nesse contexto, a integração de informações com o uso de Inteligência Artificial Generativa, especialmente por meio de Modelos de Linguagem de Grande Escala (LLMs), surge como uma alternativa promissora para orientar os usuários de forma personalizada e didática. Este trabalho tem como objetivo projetar e desenvolver uma aplicação web de gestão financeira acompanhada de uma LLM, viabilizando sua aplicação prática no contexto da educação financeira. \\
  Palavras-chave: Inteligência Artificial; LLMs; RAG; Aplicação Web; Educação Financeira.
\end{resumo}

\begin{abstract}
  Regarding finances, Brazilians face a significant challenge, 
  reflected in a lack of financial literacy. Research conducted by that a large portion of Brazilians have little or no understanding of the subject. Although traditional financial control systems help monitor income and expenses, the absence of adequate knowledge limits effective decision-making. In this context, the integration of information using Generative Artificial Intelligence, especially through Large Language Models (LLMs), emerges as a promising alternative to guide users in a personalized and didactic way. This work aims to design and develop a web-based financial management application integrated with an LLM, enabling its practical use in financial education.
\end{abstract}

\section{INTRODUÇÃO}

\subsection{MOTIVAÇÃO E JUSTIFICATIVA}

  De acordo com a pesquisa Observatório Febraban, realizada pelo Instituto de Pesquisas Sociais, Políticas e Econômicas (IPESPE) em 2025, aproximadamente 40\% dos brasileiros afirmam possuir pouco conhecimento sobre educação financeira, enquanto 15\% declaram não possuir entendimento sobre o tema, o que evidencia que a saúde financeira da população brasileira constitui um desafio estrutural relevante. Segundo a Organização para a Cooperação e Desenvolvimento Econômico \cite{oecd:2023}, baixos níveis de educação financeira impactam diretamente a capacidade dos indivíduos de realizar planejamentos econômicos adequados e tomar decisões financeiras conscientes. Nesse contexto, a pesquisa também demonstra que, embora a maioria dos entrevistados reconheça a importância da educação financeira, ainda existem limitações significativas no acesso ao conhecimento necessário para o planejamento e controle adequados das finanças pessoais \cite{febraban:2025}. \\
  Estudos demonstram que o uso de aplicações financeiras digitais apresenta-se como ferramentas que podem contribuir para a melhoria da tomada de decisões econômicas ao ampliar o acesso às informações financeiras dos usuários \cite{carlin2023mobile}. No entanto, tal solução, isoladamente, não é suficiente, uma vez que, conforme estabelece a Comissão de Valores Mobiliários \cite{cvm:2022}, a educação financeira influencia diretamente o comportamento dos indivíduos e sua capacidade de interpretar informações e tomar decisões adequadas voltadas ao bem-estar financeiro, de modo que a ausência de conhecimento financeiro pode comprometer a compreensão dos dados disponibilizados por essas ferramentas e reduzir seu potencial de contribuição para a melhoria da saúde econômica dos indivíduos. \\
  Diante desse cenário, a integração entre a centralização de informações financeiras e a utilização de Inteligência Artificial Generativa apresenta-se como uma promissora solução para auxiliar o usuário de forma mais ativa e personalizada. Segundo Li et al. (2023), Modelos de Linguagem de Grande Escala (LLMs) têm demonstrado capacidade de interpretar dados, identificar padrões e fornecer orientações em linguagem natural, o que pode contribuir significativamente para a educação financeira do usuário. \\
  Assim, surge o seguinte questionamento: como a adição de uma inteligência artificial generativa a um sistema de controle financeiro pode oferecer suporte e colaborar com a educação financeira do usuário? Mais especificamente, busca-se analisar a capacidade de um modelo de identificar problemas financeiros, fornecer orientações didáticas e contribuir para a tomada de decisões financeiras mais assertivas, conscientes e sustentáveis. Este estudo visa projetar, desenvolver e avaliar a viabilidade dessa integração em um protótipo de aplicação web focada em educação financeira personalizada, ampliando o acesso a ferramentas inteligentes de apoio à gestão econômica individual.

\subsection{OBJETIVO GERAL}

  O objetivo deste trabalho é desenvolver uma aplicação web assistida por LLMs, capaz de fornecer diagnósticos e orientações financeiras personalizadas no contexto da gestão pessoal. Busca-se analisar de que formas essa tecnologia pode interpretar informações financeiras, identificar padrões, reconhecer problemas relacionados à organização financeira e oferecer recomendações em linguagem clara e acessível. Sob esse viés, a construção do protótipo visa validar a eficácia da integração da inteligência artificial na assistência à educação financeira, demonstrando ser viável como ferramenta de apoio à gestão econômica individual, de modo a ampliar o acesso a soluções inteligentes que auxiliem os usuários na construção de hábitos financeiros conscientes, sustentáveis e alinhados aos seus objetivos pessoais.

\subsection{OBJETIVOS ESPECÍFICOS}

  \begin{itemize}
    \item Mapear os principais conceitos de finanças pessoais, estruturando uma base de conhecimento que servirá de referência para o sistema.
    \item Desenvolver a arquitetura e a interface da aplicação web capaz de coletar dados financeiros do usuário, rastrear seus hábitos e processar essas informações por meio de dashboards informativos.
    \item Integrar um modelo de LLM ao sistema por meio da técnica de Retrieval-Augmented Generation (RAG), conectando as IAs à base de dados previamente estruturada.
    \item Analisar o desempenho do modelo integrado a partir de cenários financeiros simulados, avaliando, por meio de métricas técnicas, como o tempo de resposta, e de métricas qualitativas, como a ausência de alucinações e a utilidade das recomendações geradas.
  \end{itemize}

\subsection{ESTRUTURA DO TRABALHO}

  O artigo está estruturado da seguinte forma: a seção 1 apresenta a introdução, contemplando a contextualização do tema, a justificativa/motivação da pesquisa, além dos objetivos gerais e específicos da pesquisa. A seção 2 explica a metodologia adotada, detalhando os procedimentos para o desenvolvimento do estudo, incluindo a construção da base de conhecimento em finanças pessoais e a integração da LLM por meio de RAG. A seção 3 aborda a fundamentação teórica, reunindo os principais conceitos relacionados ao estudo, como conceitos relacionados à educação financeira e às tecnologias de desenvolvimento de software e de Inteligência Artificial utilizadas no decorrer do estudo.



\section{METODOLOGIA} \label{sec:metodologia}

  Essa seção especifica os caminhos seguidos durante o processo e classifica cientificamente o estudo, as definições de arquitetura e os critérios comparativos entre as LLMs

  \subsection{CLASSIFICAÇÃO DA PESQUISA}

    Esta pesquisa possui por objetivo ser aplicada, exploratória, bibliográfica e experimental. Gil (2022) afirma que a pesquisa aplicada é voltada à aquisição de conhecimento com foco em aplicação numa situação específica. Neste escopo, a finalidade prática reside no desenvolvimento de uma solução tecnológica em formato de aplicação web capaz de sanar dúvidas e contribuir para o letramento financeiro pessoal. Além disso, a natureza exploratória é justificada pela pesquisa ter como propósito “proporcionar maior familiaridade com o problema, com vistas a torná-lo mais explícito e construir hipóteses” (GIL, 2022, p. 27). \\
    Em relação à fundamentação teórica para o mapeamento dos conceitos de finanças e economia comportamental, será adotada a pesquisa bibliográfica, que “é elaborada com base em material já publicado” (GIL, 2022, p. 29). De forma complementar, adota-se a pesquisa experimental que, de acordo com Gil (2022), tem como propósito determinar um objeto de estudo, selecionar as variáveis que o influenciam e definir formas de controle, presentes por meio da prototipagem do software, em que, dentro de cenários controlados, validam-se as interações entre o usuário, o sistema e o modelo de inteligência artificial generativa

  \subsection{ARQUITETURA DO SISTEMA}

    Para tornar o teste dos modelos viável em um cenário em que um usuário do sistema fornece as informações, será necessário o desenvolvimento de um protótipo fundamentado em uma arquitetura cliente-servidor. \\
    A interface de usuário adotará a tecnologia do ecossistema React, por meio do framework Next.js e TypeScript. Essa escolha garante a tipagem estática do código e de componentes interativos na renderização dos dashboards financeiros. \\
    A lógica de negócios e a comunicação com o banco de dados PostgreSQL serão estruturadas em uma API Python, utilizando Django REST. O servidor é responsável por receber os dados inseridos pelo usuário, processá-los, estruturá-los em um prompt em formato JSON e gerenciar as requisições HTTP para o modelo de linguagem.

  \subsection{INTEGRAÇÃO DA INTELIGÊNCIA ARTIFICIAL E RAG}

    A inteligência artificial será integrada utilizando a técnica de Retrieval-Augmented Generation (RAG), utilizando uma abordagem baseada em recuperação contextual de informações. Para a implementação da solução, será adotada uma arquitetura RAG 2.0, na qual os dados recuperados passam por uma etapa adicional de processamento e contextualização antes de serem enviados ao modelo de linguagem, permitindo respostas mais específicas e alinhadas ao domínio da aplicação. \\
    A base de conhecimento será estruturada utilizando um banco de dados vetorial, responsável pelo armazenamento de embeddings gerados a partir de informações financeiras do usuário e conteúdos técnicos relacionados a educação financeira. Entre as tecnologias consideradas para essa finalidade estão ferramentas como FAISS (Facebook Al Similarity Search) ou ChromaDB, que possibilitam a indexação e recuperação eficiente de informações semânticas relevantes. \\
    O fluxo de integração entre os componentes da arquitetura será implementado utilizando o framework LangChain, responsável pela orquestração das etapas de geração de embeddings, recuperação de documentos relevantes e construção do contexto enviado ao modelo de linguagem. \\
    O modelo de linguagem será integrado por meio da API da OpenAI, sendo responsável pela interpretação dos dados recuperados e pela geração das respostas em linguagem natural. Dessa forma, o sistema anexará ao contexto da requisição tanto o histórico financeiro do usuário quanto informações provenientes da base de conhecimento, permitindo que o modelo utilize dados específicos durante o processo de geração. Essa estratégia contribui para reduzir riscos de alucinação e aumentar a consistência das respostas produzidas.


\section{FUNDAMENTAÇÃO TEÓRICA}

Esta seção do estudo apresenta os principais conceitos e tecnologias que fundamentam o desenvolvimento deste trabalho. Inicialmente, são abordados aspectos da Engenharia de Software, seguidos por conceitos da área de Economia com foco na saúde financeira pessoal. Por fim, são apresentados os fundamentos de Inteligência Artificial com ênfase em LLMs e em técnicas como Retrieval-Augmented Generation (RAG).

  \subsection{ECONOMIA}
  Conforme aponta Marchevsky (2023), as ideias clássicas de Adam Smith sobre o livre mercado e a divisão do trabalho ainda são o ponto de partida para entender como a economia gera valor. No entanto, o autor mostra que, hoje em dia, a economia vai muito além das grandes indústrias e da produção nacional, passando a olhar de perto para a vida financeira das pessoas comuns no seu cotidiano. Nesse viés, o autor acredita que nos tempos contemporâneos a realidade exige que cada indivíduo saiba como equilibrar seu próprio bolso, controlando o consumo e desenhando um planejamento financeiro realista, pois é exatamente no dia a dia que variáveis como juros, inflação e o uso do crédito mexem com o orçamento das famílias. Por isso, o estudo conclui que entender como o mercado funciona na prática virou uma questão de sobrevivência: quando as pessoas não dominam esses conceitos, as chances de cair na armadilha do endividamento e da inadimplência disparam, o que acaba prejudicando a saúde financeira dos indivíduos.

    \subsubsection{JUROS}

      Na definição de Assef (2024), os juros representam o custo do dinheiro ao longo do tempo, funcionando como a remuneração pelo uso de um capital em operações de empréstimo, financiamento ou investimento. O autor explica que esse custo pode ser calculado de forma simples ou composta, impactando diretamente o valor final que será pago ou recebido pelo investidor. Complementando essa visão, Kobori (2024) destaca em sua obra que os juros exercem um papel essencial na dinâmica e nas decisões financeiras individuais. Para Kobori, as taxas de juros influenciam diretamente o poder de consumo das famílias, o acesso ao crédito e, principalmente, a capacidade de formação de patrimônio da população ao longo do tempo.

    \subsubsection{INFLAÇÃO}

      Conforme dados oficiais do Banco Central do Brasil (2024), a inflação é caracterizada pelo aumento generalizado e contínuo dos preços de bens e serviços em uma economia ao longo do tempo. Para a autarquia, esse fenômeno reduz de forma direta o poder de compra da moeda, fazendo com que uma mesma quantia adquira menos produtos. Em consonância com essa visão, Kobori (2024) explica em sua obra que a inflação está diretamente associada ao desequilíbrio entre a oferta e a demanda monetária. O autor ressalta que a expansão da base monetária, quando não é acompanhada pelo crescimento real da produção, tende a intensificar a elevação dos preços. Esse cenário é um reflexo direto das decisões de política monetária no Brasil, cujos impactos no controle inflacionário e na estabilidade econômica são fundamentais para proteger a renda da população (FLORA; SANTOLIN, 2024).

    \subsubsection{INADIMPLÊNCIA}

      Na perspectiva de Kobori (2024), a inadimplência é o reflexo direto do descumprimento de obrigações financeiras por parte de indivíduos ou empresas, identificada pelo atraso ou não pagamento de dívidas no prazo acordado. Em sua obra, o autor explica que esse fenômeno está frequentemente associado à falta de planejamento pessoal e ao uso inadequado de linhas de crédito pelas famílias. Sob essa ótica, Kobori defende que a educação financeira exerce um papel fundamental e preventivo, funcionando como a principal ferramenta para mitigar o endividamento excessivo e garantir a sustentabilidade da saúde financeira da população.

    \subsubsection{SAÚDE FINANCEIRA PESSOAL}

      Na visão de Lazaretti, Castanheira e Maschio (2024), a saúde financeira pessoal baseia-se diretamente na capacidade que o indivíduo tem de equilibrar sua renda com as despesas diárias, mantendo o controle sobre suas obrigações financeiras imediatas e futuras. Para os autores, esse conceito está intimamente ligado à organização, ao uso consciente do crédito e à tomada de decisões que garantam estabilidade ao longo do tempo. Eles apontam que o desenvolvimento de uma saúde financeira robusta depende de programas contínuos de alfabetização financeira, os quais são fundamentais para capacitar o cidadão a mitigar os riscos do endividamento e a exercer escolhas de consumo muito mais conscientes e sustentáveis.

    \subsubsection{ECONOMIA COMPORTAMENTAL}

      De acordo com Housel (2021), o sucesso financeiro está muito mais ligado ao comportamento humano e à maneira como lidamos com o dinheiro do que ao conhecimento técnico ou matemático em si. Sob essa ótica, a Economia Comportamental ganha relevância ao estudar como fatores psicológicos, emocionais e cognitivos influenciam as escolhas financeiras dos indivíduos. O autor argumenta que as decisões do dia a dia não são tomadas puramente com base em planilhas ou fórmulas, mas são moldadas diretamente pelas emoções, pelos hábitos cotidianos e pelas experiências pessoais de cada investidor.

  \subsection{ENGENHARIA DE SOFTWARE}

    De acordo com Ciriano e Scombatti Pinto (2025), o avanço dos sistemas computacionais e a crescente integração entre tecnologias ampliaram a necessidade de métodos estruturados para o desenvolvimento de software. Os autores retomam fundamentos clássicos da área ao destacar que Pressman e Maxim (2020) definem a Engenharia de Software como a aplicação sistemática de princípios, métodos e práticas voltadas ao desenvolvimento e manutenção de sistemas computacionais de forma organizada e eficiente. Dessa forma, a Engenharia de Software assume papel essencial na construção de aplicações modernas ao fornecer mecanismos que favorecem organização estrutural, comunicação entre componentes e maior flexibilidade no desenvolvimento de sistemas.

    \subsubsection{ARQUITETURA CLIENTE-SERVIDOR}

      Ciriano e Scombatti Pinto (2025) apontam que a evolução das arquiteturas de software ocorreu em resposta às demandas por escalabilidade, flexibilidade e melhor organização estrutural dos sistemas computacionais. Os autores ressaltam que aplicações modernas passaram a exigir modelos capazes de distribuir responsabilidades e facilitar a comunicação entre componentes, tornando a arquitetura cliente-servidor uma das abordagens mais empregadas em ambientes distribuídos. \\
      Nessa perspectiva, os autores destacam que estudos clássicos de Sommerville (2019) definem a arquitetura cliente-servidor como um modelo em que as responsabilidades do sistema são divididas entre clientes e servidores. Enquanto o cliente é responsável pela interação com o usuário e pela solicitação de recursos, o servidor realiza o processamento das requisições, executa regras de negócio e disponibiliza os serviços solicitados. Essa separação estrutural reduz a dependência entre componentes e contribui para uma organização mais eficiente do sistema. \\
      Complementando essa visão, os autores (2025) também reforçam, a partir dos fundamentos apresentados por Bass, Clements e Kazman, que modelos arquiteturais modernos buscam favorecer características como modularidade, facilidade de manutenção e escalabilidade. Sob essa perspectiva, aplicações web utilizam amplamente a arquitetura cliente-servidor para separar a camada responsável pela interface do usuário das camadas relacionadas ao processamento de dados e gerenciamento das informações, permitindo o desenvolvimento de aplicações mais flexíveis e adequadas a ambientes distribuídos

    \subsubsection{API’S E INTEGRAÇÃO DE SISTEMAS}

      Para Ciriano e Scombatti Pinto (2025), a evolução das arquiteturas modernas ampliou a necessidade de mecanismos que permitam comunicação eficiente entre diferentes componentes e serviços computacionais. Nesse cenário, a integração entre sistemas passou a desempenhar papel importante para garantir interoperabilidade e flexibilidade em aplicações distribuídas.
      Sob essa perspectiva, estudos apresentados por Sommerville (2019) definem APIs (Application Programming Interfaces) como mecanismos responsáveis por estabelecer regras e padrões de comunicação entre sistemas distintos, permitindo o compartilhamento padronizado de dados e funcionalidades. Essa abordagem reduz o acoplamento entre aplicações e facilita a integração entre diferentes serviços, contribuindo para o desenvolvimento de sistemas mais organizados e escaláveis.

  \subsection{INTELIGÊNCIA ARTIFICIAL}
  
    A gênese da Inteligência Artificial (IA) remonta à década de 1960 com o famoso Teste de Turing. Esse modelo funcionou como um ponto de partida para discutir se as máquinas poderiam, de fato, simular o pensamento humano (RUSSELL; NORVIG, 2022). No início, cientistas como John McCarthy acreditavam que, para uma máquina ser inteligente, ela precisaria seguir regras lógicas rígidas e tentar imitar o senso comum humano. Porém, com o passar dos anos, esse modelo se mostrou limitado (GABRIEL FILHO, 2023). Hoje, a IA moderna funciona de um jeito totalmente diferente. Em vez de apenas seguir ordens pré-programadas em um banco de dados fixo, os sistemas atuais utilizam o Aprendizado de Máquina para analisar grandes volumes de dados e aprender com eles de forma autônoma, adaptando-se a novas situações (FACELI et al., 2025).

    \subsubsection{MACHINE LEARNING}

      Conforme destacado por Barbierato e Gatti (2024), o Machine Learning (Aprendizado de Máquina) é uma subárea da Inteligência Artificial voltada ao desenvolvimento de algoritmos capazes de aprender a partir de dados e generalizar comportamentos para novas situações por meio de métodos estatísticos e computacionais. Por meio da identificação de padrões e relações, esses sistemas conseguem realizar previsões, reconhecer comportamentos e auxiliar processos de tomada de decisão em diferentes contextos. Nesse sentido, os avanços recentes em aprendizado de máquina têm impulsionado significativamente a evolução dos sistemas de Inteligência Artificial e ampliado sua aplicação em diversas áreas da sociedade.

    \subsubsection{LLM}

      Segundo Minaee et al. (2024), os Modelos de Linguagem de Grande Escala (LLMs) são sistemas avançados de Inteligência Artificial treinados com grandes volumes de dados, capazes de compreender, processar e gerar linguagem natural de forma semelhante à humana. Além disso, estudos recentes destacam a evolução desses modelos para arquiteturas multimodais, capazes de integrar diferentes tipos de dados, como texto, imagens e áudio, ampliando suas capacidades de interpretação e interação contextual (HUANG et al., 2024). Conforme os autores, os modelos utilizam técnicas de aprendizado profundo para identificar padrões linguísticos e contextuais, permitindo sua aplicação em tarefas como geração de texto, análise de informações e suporte à tomada de decisão.

    \subsubsection{RAG}

      Segundo Gao et al. (2024), a técnica Retrieval-Augmented Generation (RAG) consiste em um processo de otimização das respostas geradas por LLMs, no qual o modelo consulta uma base de dados externa e confiável durante a geração das respostas. Dessa forma, o artigo conceitua que o RAG combina os conhecimentos pré-treinados do modelo com informações atualizadas e específicas, reduzindo erros e aumentando a qualidade das respostas. Para os autores, essa abordagem melhora a precisão e a confiabilidade dos sistemas baseados em linguagem natural.

\nocite{*}
\bibliographystyle{sbc}
\bibliography{sbc-template}

\end{document}
